\section{三个扩展接口及其附加条件}
\label{app:extensions}

本附录研究固定表示下的三个扩展接口：标准 Borel 可测因子的精确闭合、随机
通道的一步边缘交织，以及有限状态连续时间的生成元与半群关系。相应的系统族
$\Gamma_N$ 判据还需要各自的复杂度、紧性、最优达到与可检验性条件。正文主
定理适用于有限状态、离散时间和确定性任务保真商。标准 Borel kernel 与正则
条件概率的背景可参见\citet{Kallenberg2021}；标准 Borel 空间和 Borel
等价关系的描述集合论背景见\citet{Kechris1995}。

\subsection{给定标准 Borel 因子的精确闭合}

\begin{proposition}[标准 Borel 空间中的固定因子判据]
\label{prop:borel-fixed-factor}
设 $X,Y$ 为标准 Borel 空间，$s:X\to Y$ 为可测满射，
$K:X\rightsquigarrow X$ 为 Markov kernel。对每个 Borel 集
$A\subseteq Y$，假设
\[
K(x,s^{-1}(A))
\]
在 $s$ 的每条纤维上常值，并且由此诱导的函数
\[
y\longmapsto K(x,s^{-1}(A)),
\qquad s(x)=y,
\]
是可测的。则
\[
L(y,A):=K(x,s^{-1}(A)),
\qquad s(x)=y,
\]
良定义且给出唯一 Markov kernel $L:Y\rightsquigarrow Y$，满足
\begin{equation}
KQ_s=Q_sL.
\label{eq:borel-fixed-factor}
\end{equation}
反之，式~\eqref{eq:borel-fixed-factor}蕴含上述纤维常值性与诱导函数的
可测性。
\end{proposition}

\begin{proof}
纤维常值性保证 $L$ 不依赖代表元。固定 $y$ 后，
$A\mapsto L(y,A)$ 从 $K(x,\cdot)$ 经逆像 $s^{-1}$ 得到，因而是概率测度；
对每个 $A$ 的可测性由假设给出，所以 $L$ 是 Markov kernel。满射性给出
唯一性，交换式按每个 Borel 集逐点验证。反向结论直接读取交换式两侧在
$x$ 处对事件 $A$ 的取值即可。
\end{proof}

该判据把可测性作为显式条件。一般 Borel 等价关系可以是 nonsmooth 的，此时
抽象商缺少实现上述命题所需的标准 Borel 结构\citep{Kechris1995}；逐纤维
定义的值组成 kernel 还需要整体可测性。在观测语言中，给定因子对应子
$\sigma$-代数 $s^{-1}(\mathcal B(Y))$；精确闭合要求 Markov 算子保持其中的
有界可测观测量。标准 Borel 商的无损实现以及商空间表述与子 $\sigma$-代数
表述的等价，需要可数生成等额外正则性。

命题~\ref{prop:borel-fixed-factor}给出\emph{固定 $s$} 的精确因子判据。
一般 Borel 系统族判据还需指定与有限块数对应的复杂度泛函、紧致候选因子类、
目标函数下半连续性、最优值达到及可终止检验程序。

\subsection{随机粗粒化的一步交织接口}

令 $Q:X\rightsquigarrow Y$ 本身为 Markov kernel。随机粗粒化的精确交换式为
\begin{equation}
KQ=QL.
\label{eq:random-intertwining}
\end{equation}
它表示：先微观演化再随机编码，与先随机编码再按 $L$ 演化，具有相同的一步
输出边缘。Rogers--Pitman 型 Markov functions 提供了与 link kernel
相容的投影框架\citep{RogersPitman1981}；随机映射及其复合也可由 Markov
category 组织\citep{Fritz2020}。

随机 $Q$ 的完整判据需要处理以下四项结构：
\begin{enumerate}
  \item 随机通道通常对应软编码，其复杂度需要块数以外的泛函；
  \item 任务保真、信息约束或非退化条件用于排除常值噪声通道产生的任务无关
  平凡交织：
  当 $Q(x,\cdot)=\nu$ 时，只要 $\nu L=\nu$，式~\eqref{eq:random-intertwining}
  就成立；
  \item 交换式的解集可以包含多个 $L$，上述常值通道已经给出例子；
  \item 过程级结论需要微观路径与宏观路径的联合耦合，并需刻画哪些宏观初始
  分布可由微观初态实现。
\end{enumerate}

近似情形仍可使用 $d_\infty(KQ,QL)$。固定 $Q$ 且 $Y$ 有限时，最优 $L$
仍由有限概率单纯形的紧致性达到；同时优化 $Q$ 的完整判据还必须指定
任务通道保真条件、随机表示的复杂度或信息约束、排除平凡通道的非退化条件、
候选 $Q$ 的紧致性，以及需要路径结论时的 link consistency 或联合耦合。
这些选择共同确定随机通道版本 $\Gamma_N$ 的具体定义。

\subsection{有限状态连续时间的固定因子关系}

令 $G$ 与 $\bar G$ 分别是有限状态空间 $X,Y$ 上的保守生成元，$Q$ 为任意
类型匹配的 Markov 通道。对行和为零的矩阵 $A$，记
\[
\|A\|_{\infty,\TV}
:=\max_x\frac12\sum_y|A(x,y)|.
\]

\begin{proposition}[生成元交织、半群交织与有限时间误差]
\label{prop:continuous-time-fixed-factor}
若
\begin{equation}
GQ=Q\bar G,
\label{eq:generator-intertwining}
\end{equation}
则对所有 $t\ge0$，
\begin{equation}
e^{tG}Q=Qe^{t\bar G}.
\label{eq:semigroup-intertwining}
\end{equation}
反之，若式~\eqref{eq:semigroup-intertwining}在 $t=0$ 的某个右邻域成立，则
式~\eqref{eq:generator-intertwining}成立。更一般地，若
\[
\|GQ-Q\bar G\|_{\infty,\TV}\le\eps,
\]
则对每个 $t\ge0$，
\begin{equation}
\|e^{tG}Q-Qe^{t\bar G}\|_{\infty,\TV}
\le t\eps.
\label{eq:continuous-time-linear-bound}
\end{equation}
\end{proposition}

\begin{proof}
精确方向可由矩阵幂级数逐项推出；反向在 $t=0$ 取右导数即可。近似方向使用
有限维 Duhamel 公式
\[
e^{tG}Q-Qe^{t\bar G}
=\int_0^t e^{(t-r)G}(GQ-Q\bar G)e^{r\bar G}\dd r.
\]
左右两个 Markov 半群分别对最大行 TV 范数和行和为零的有符号测度非扩张，
故被积式范数至多为 $\eps$；积分得到式~\eqref{eq:continuous-time-linear-bound}。
\end{proof}

这个命题适用于固定通道和固定有限时间。连续时间系统族还需定义压缩复杂度与
随规模增长的统一时间界；一般状态空间进一步涉及生成元定义域、强连续性、
闭包和 martingale problem，可参见\citet{EthierKurtz1986}。这些解析条件
构成连续时间扩展的附加前提。
